\begin{center}

{\footnotesize
\renewcommand{\arraystretch}{1.15}
\setlength{\tabcolsep}{4pt}

\begin{longtable}{
|>{\raggedright\arraybackslash}p{8.80cm}
|>{\centering\arraybackslash}p{2.80cm}
|>{\centering\arraybackslash}p{2.00cm}|
}

\caption{Ringkasan evaluasi RM3 pada penulisan graf, pembaruan
\textit{lifecycle}, retrieval, dan penggunaan memori pada jawaban}
\label{tab:v2-memory-result}
\\
\hline

\rowcolor{tableheadgray}
\textbf{Metrik}
&
\textbf{Kondisi}
&
\textbf{Nilai}
\\
\hline
\endfirsthead

\multicolumn{3}{c}{
\textit{Lanjutan Tabel \ref{tab:v2-memory-result}}
}
\\
\hline

\rowcolor{tableheadgray}
\textbf{Metrik}
&
\textbf{Kondisi}
&
\textbf{Nilai}
\\
\hline
\endhead

\hline
\multicolumn{3}{r}{
\textit{Bersambung ke halaman berikutnya}
}
\\
\endfoot

\hline
\endlastfoot

\multicolumn{3}{|l|}{
\textit{Penulisan graf}, $n = 12$ skenario
}
\\
\hline

F1 node
&
Tidak berlaku
&
0,620
\\
\hline

Presisi node
&
Tidak berlaku
&
0,556
\\
\hline

Recall node
&
Tidak berlaku
&
0,702
\\
\hline

F1 relasi
&
Tidak berlaku
&
0,361
\\
\hline

Akurasi label \texttt{subject\_role}
&
Tidak berlaku
&
85,7\%
\\
\hline

Akurasi label \texttt{trigger\_category}
&
Tidak berlaku
&
100,0\%
\\
\hline

\multicolumn{3}{|l|}{
\textit{Pembaruan lifecycle}, $n = 12$ skenario
}
\\
\hline

Ketepatan tindakan pembaruan
&
Tidak berlaku
&
33,3\%
\\
\hline

Ketepatan konten baru
&
Tidak berlaku
&
50,0\%
\\
\hline

Konten lama berhasil dinonaktifkan
&
Tidak berlaku
&
25,0\%
\\
\hline

\multicolumn{3}{|l|}{
\textit{Retrieval}, $n = 42$ kueri pada dua kondisi
}
\\
\hline

P@5
&
Hibrid
&
0,073
\\
\hline

P@5
&
Vektor saja
&
0,000
\\
\hline

Recall@5
&
Hibrid
&
0,267
\\
\hline

Recall@5
&
Vektor saja
&
0,000
\\
\hline

MRR
&
Hibrid
&
0,267
\\
\hline

MRR
&
Vektor saja
&
0,000
\\
\hline

nDCG@5
&
Hibrid
&
0,257
\\
\hline

nDCG@5
&
Vektor saja
&
0,000
\\
\hline

Tingkat memori usang
&
Hibrid
&
0,0\%
\\
\hline

Keberhasilan kueri relasional atau kausal
&
Hibrid
&
29,2\%
\\
\hline

Keberhasilan abstensi
&
Hibrid
&
45,2\%
\\
\hline

\multicolumn{3}{|l|}{
\textit{Penggunaan memori pada jawaban},
$n = 24$ kasus pada dua kondisi
}
\\
\hline

Jawaban berdasar memori
&
Hibrid
&
95,8\%
\\
\hline

Jawaban berdasar memori
&
Vektor saja
&
95,8\%
\\
\hline

Tingkat kontradiksi
&
Hibrid
&
8,3\%
\\
\hline

Tingkat kontradiksi
&
Vektor saja
&
4,2\%
\\
\hline

Keberhasilan relasional atau kausal
&
Hibrid
&
66,7\%
\\
\hline

Keberhasilan relasional atau kausal
&
Vektor saja
&
62,5\%
\\
\hline

Personalisasi tanpa dasar
&
Hibrid
&
29,2\%
\\
\hline

Personalisasi tanpa dasar
&
Vektor saja
&
29,2\%
\\
\hline

Tingkat keyakinan usang
&
Hibrid
&
8,3\%
\\
\hline

Tingkat keyakinan usang
&
Vektor saja
&
8,3\%
\\
\hline

Abstensi benar
&
Hibrid
&
25,0\%
\\
\hline

Abstensi benar
&
Vektor saja
&
25,0\%
\\
\hline

\end{longtable}

\begin{minipage}{\textwidth}
\footnotesize
\textit{Catatan:}
Kondisi tidak berlaku digunakan pada blok evaluasi yang tidak
membandingkan retrieval hibrid dengan retrieval vektor saja.
Nilai yang lebih tinggi menunjukkan hasil yang lebih baik, kecuali
untuk tingkat kontradiksi, tingkat memori usang, personalisasi tanpa
dasar, dan tingkat keyakinan usang.
\end{minipage}
}

\end{center}